\uuid{wT5O}
\titre{Strucutre de sous-espace vectoriel}
\auteur{hanani ; mzari}
\organisation{exo7}
\licence{CC BY-NC-SA 4.0}
\datecreate{2020-01-20}
\niveau{L1}
\module{Algèbre}
\chapitre{Espace vectoriel}
\sousChapitre{Définition, sous-espace}
\difficulte{1}
\theme{Espaces vectoriels, Strucutre de sous-espace vectoriel}
\exo7qcmid{lille-2.espaces-vectoriels.001}

\contenu{
\begin{question}
\qclassification{106.01}
\qtags{motcle=Strucutre de sous-espace vectoriel}
Soit $E=\{(x,y)\in \Rr^2 \, ; \;  x+y=1\}$, muni des opérations usuelles. Quelles sont les assertions vraies ?
\begin{answers}
\bad{$E$ est un espace vectoriel, car  $E$ est un sous-ensemble de l'espace vectoriel $\Rr^2$.}
\good{$E$ n'est pas un espace vectoriel, car $(0,0)\notin E$.}
\good{$E$ n'est pas un espace vectoriel, car $(1,0)\in E$, mais $(-1,0)\notin E$.}
\good{$E$ n'est pas un espace vectoriel, car $(1,0)\in E$ et $(0,1)\in E$, mais $(1,1)\notin E$.}
\end{answers}
\begin{explanations}
$E$ n'est pas un sous-espace vectoriel de $\Rr^2$, puisque $(0,0) \notin E$.\\
$E$ n'est pas stable par multiplication par un scalaire : $(1,0) \in E$ , mais,  $-(1,0)=(-1,0) \notin E$.\\
$E$ n'est pas stable par addition : $(1,0)  \in E $ et $(0,1)  \in E $, mais $(1,0)+(0,1)=(1,1) \notin E$.
\end{explanations}
\end{question}
}
